\chapter{120}\label{section-120}

Stazione Genua Nervi.

Der Bahnhof lag direkt an der Meeresküste.

Louis bog neben dem Gebäude auf einen langen Parkplatz ein. Er war fast
leer. Es war offensichtlich. Der diesjährige Sommer war Geschichte. Nur
wenige Touristen schlenderten auf dem Platz herum.

Louis liess das Seitenfenster hinuntergleiten. Die Aussenluft drang
angenehm kühl in den Wagen.

«Uff.»

Er legte sich die Hände in den Schoss.

«Machen wir’s kurz, Louis. Ich checke im Hotel ein und warte dort auf
dich,» Joh fingerte einen kleinen Zettel aus der Brusttasche seines
Hemds, «für dich, Manuela hat uns die Adresse des Hotels
aufgeschrieben.»

Er reichte Louis die Notiz, öffnete Google-Map auf dem Handy und gab die
Angaben ein.

«Also, \emph{Hotel Villa Bonera, Via Sarfatti 8,} hier, schau,» Joh
drehte das Display zu Louis, «das sind nur ungefähr 20 Gehminuten vom
\emph{Torre Gropallo} entfernt.»

«Danke. Und hier,» Louis beugte sich vor und führte seinen Zeigefinger
über dem Handy dem Weg entlang, «hier ist gleich der Wachturm gegen
Osten.»

Joh nickte.

«Und wie vereinbart, Louis, ich bleibe im Hotel. Sobald du mich brauchen
solltest, rufst du an. Und sonst, vergiss mich einfach. Ich kann warten,
auch wenn du Stunden brauchst, ok?»

Louis nickte abermals und Joh stieg aus. Er ging um den Wagen herum und
stellte sich vor die Führertür. Louis schälte sich aus dem Sitz und Joh
zog ihn mit einem Ruck zu sich hinaus. Er fasste ihn bei den Schultern,
zog ihn an sich und löste sich auch gleich wieder.

«Ich glaube an dich,» und mit den Worten setzte sich Joh ans Steuer und
schloss die Wagentür. Dann bog er zügig auf die \emph{Viale delle Palme}
ein und fuhr Richtung Innenland davon.

\emph{Nichts denken}, murmelte Louis vor sich hin\emph{, nichts denken,
einfach gehen.} Während seiner ersten Schritte gelang es. Bald aber
begannen diverse Stimmen lästig um seine Aufmerksamkeit zu streiten. Die
Letzte trat wie eine willkommene Trostspenderin auf. \emph{Joh ist da.
Du kannst ihn jederzeit rufen.}

Der Einstieg in die kleine Unterführung fand Louis auf Anhieb. Er zielte
als geschwungener Weg abwärts und direkt auf die \emph{Passeggiata}
zu. Kaum war sein Blick nach unten frei, tat sich vor ihm ein
rechteckiges, steinernes Panoramafenster zum Meer auf.

Louis liess seinen Blick über die Promenade schweifen, erst nach rechts,
dann nach links. Da war er, dieser Jahrhunderte alte Turm. Louis
schauderte. Dort also würde er Giorgia treffen.

Die Länge des Spazierweges übertraf seine Vorstellung. Vor ihm lag das
Meer wie ein Silberteppich, bewegte sich gelassen auf und ab, klatschte
unregelmässig kleine Wassermassen an die Felsen, die die Wellen gleich
wieder verschluckten.

Louis stützte sich auf dem Geländer ab. Er krallte die Finger um das
Eisenrohr und nahm ein Zittern seines Körpers wahr. Vielleicht lag es
einzig an der Besonderheit dieses Ortes, oder, dass er ganz alleine
dastand, kein Mensch weit und breit zu sehen war, oder daran, dass er
einfach erschöpft war. Er weinte hemmungslos. Dabei blieb sein Blick auf
dem unweit vor ihm stehenden Turm liegen.

Er suchte nach einem Taschentuch, als sich über ihm der Himmel aufriss.
Erst waren es nur einzelne spitze Sonnenstrahlen. Als würden sie gerade
frei gelassen, peilten sie blitzschnell die Meeresoberfläche an. Dann
sah er zu, wie sich das dunkle Meer mit einer Glanzschicht überzog.
Fasziniert liess er seinen Blick über die unermesslich weite Oberfläche
gleiten.

Louis löste sich vom Geländer, trat zurück und atmete laut aus. Der
kurze Moment ganz mit sich alleine war wohltuend. Er strich sich das
Haar aus dem Gesicht und hob seinen Rucksack hoch. Dann streckte er sich
durch.

Er lag gut in der Zeit.

Mit festen Schritten ging er los. Von vorne kam ihm eine Menschengruppe
entgegen. Sie verzettelte sich, kam wieder zusammen und bald konnte er
mehrere Teenager wild durcheinander grölen und herumalbern hören. Ein
Mädchen kam ihm entgegengerannt. In der einen Hand schwang sie eine
Schirmmütze durch die Luft.

«Hol sie dir doch, du Pfeife.»

Ein Junge schrie ihr eine Ladung Schimpfwörter hinterher. Auf Louis’
Höhe angekommen, versuchte sie sich hinter ihm zu verstecken. Er hörte
sie heftig atmen. Schliesslich hatten die andern sie eingeholt und
hinter Louis fand die Szene ihre Fortsetzung.

Louis grinste, ach, Kinder, wie habt ihr es gut.

Nach einigen Metern blieb er stehen. Vor dem Bauwerk war halbrund ein
Platz angelegt. Sein Herz raste. Giorgia. Es war soweit. Nach drei
Monaten, nach katastrophalen Wochen, für sie wie für ihn, sahen sie sich
wieder.

Je näher er dem Turm kam, desto unsicherer wurde er und verlangsamte
sein Tempo.

Der Platz war leer.

Nervös suchte er die Umgebung ab. Keine Menschenseele weit und breit. Er
trat näher. Dann sah er jemanden um die Ecke und auf ihn zu kommen. Ein
älterer Mann schob einen Rollstuhl vor sich hin. Darin sass Giorgia.
Tatsächlich, das war sie.

Louis starrte ihnen entgegen. Als Giorgia ihm zuwinkte und dazu sogar
lächelte, wäre er ihnen am liebsten entgegengestürmt. Der Mann stoppte
den Rollstuhl schliesslich vor Louis’ Füssen und strahlte ihn an.

Giorgia streckte Louis die Hand entgegen.

«Hi, Louis, wie schön, dich zu sehen», sie zog ihn zu sich hinunter und
küsste ihn auf die Wangen. Verdattert richtete er sich wieder auf,
suchte mit offenem Mund nach Worten und sah Giorgia auf den Mann zeigen.

«Das ist \emph{Zio Alessio}\footnote{Onkel Alessio},» sie deutete auf Louis, «und das ist nun
Louis.»

Erstmal fiel Louis ein Stein vom Herzen. Giorgias Stimme wie auch ihr
Sprechtempo hatten sich nicht verändert. Sie redete höchstens leiser als
gewöhnlich.

Sie begrüssten sich. Der Mann liess Louis’ Hand vorerst nicht mehr los.

«Aber natürlich, ich erinnere mich an dich, Louis,» er schien sich zu
freuen, «dich habe ich doch an der Geburtstagsfeier in Mailand gesehen.
Das war?»

Nach wie vor hielt er Louis’ Hand. Er schien nach einer Jahreszahl zu
suchen.

Giorgia kam ihm zuvor.

«Klar, Zio, das war vor drei Jahren, als Maman sechzig wurde.»

«Ma certo\footnote{«Aber sicher.»}.»

Jetzt liess der Mann Louis` Hand los.

Louis’ Lächeln misslang. Er war tief verunsichert, weil er die
Freundlichkeit der beiden nicht einordnen und sich auch nicht an diesen
Onkel erinnern konnte. Das Wiedersehen begann zu hundert Prozent anders,
als er es sich vorgestellt hatte.

Das Gesicht des Mannes sagte ihm nichts.

«Ja, genau, jetzt erinnere ich mich wieder.»

Die Worte waren von alleine über seine Lippen gekommen. Warum hatte er
das gesagt? Giorgia sah verunsichert aus. Sie hatte ihn bestimmt
durchschaut. Louis wurde heiss.

«Zio Alessio hat zu tun, Louis. Er holt mich später wieder ab.»

Sie drehte den Kopf hoch und nach hinten.

«Danke, Zio.»

«Gern, \emph{amore}, ruf einfach an.»

Er küsste sie aufs Haar.

Dann wandte er sich an Louis, reichte ihm die Hand und zog ihn nah an
sich.

«Pass auf sie auf, Louis, bis später.»

Sie waren alleine. Jetzt erst schaute Louis sie wirklich an. Wie schön
sie war. Ihre schwarzen Locken waren zu einem Dutt hochgesteckt. Ein
dezent gemalter Lidstrich betonte ihre dunklen Augen und auf ihren
Lippen lag ein rötlicher Glanz. Sie trug den hellen Mantel, den er immer
an ihr gemocht hatte. Er konnte kaum mehr wegsehen. Wie oft er sie sich
über die Wochen vorgestellt und sich gefürchtet hatte, einem Wrack zu
begegnen, einer gebrochenen Frau. Nichts davon. Höchstens blasser sah
sie aus und ihre schlanken Beine wirkten dünner.

Giorgias Gesichtsausdruck hatte sich nach der Verabschiedung ihres
Onkels verändert. Ja, Joh hatte richtig vermutet, auch sie war unsicher.

Giorgia streckte ihre Hand nach ihm aus. Louis legte seine hinein.

«Ich bin jetzt Rollstuhlfahrerin,» sie versuchte, neutral zu klingen,
«und für heute kannst du mein Fahrer sein, ok?»

«Bin ich gern. Hilf mir, wenn ich etwas falsch mache.»

Sie lösten die Hände voneinander.

«Wohin soll ich dich schieben? Du kennst dich bestimmt bereits gut aus,»
Louis ging in die Knie, um auf Augenhöhe weiterzusprechen, «deine Mam
hat mir erzählt, dass du bei deinem Onkel wohnst und hier die Reha
weiterführst?»

«Ja, Zio Alessio und seine Freundin haben mich für drei Monate hier nach
Nervi eingeladen. Wegen des Klimas, meinten sie, aber natürlich in
erster Linie, damit mich Dr. Trussardi wieder fit macht. Er ist mit
Alessio befreundet und bekannt für seine Reha-Erfolge.»

«Das hört sich gut an.»

Sie nickte. Der Ausdruck in ihren Augen hatte nicht wirklich zu ihrer
Sprechmelodie gepasst. Louis stand auf.

«Was hältst du davon, wenn du mich erstmal über die Promenade rollst und
wir uns dabei voneinander erzählen?»

«Eine gute Idee. Hast du warm genug, ich meine, du sitzt ja vor allem
im, ehm,» es fiel Louis schwer, das Wort auszusprechen.

«...im Rollstuhl, Louis. Mach dir nicht zu viele Gedanken. Ich hatte
viel Zeit, mich mit dem Ding anzufreunden. Ich sag dir,» und ganz
plötzlich glänzten ihre Augen in dieser vertrauten Schönheit, «in drei,
vier Monaten werde ich wieder gehen können und den Stuhl im Meer
versenken.»

Der Trotz stand ihr gut.

«Das traue ich dir zu, Giorgia, so gut kenn ich dich.»

Louis nahm ein feines Zucken auf ihrem Gesicht wahr.

Giorgia löste die Bremse des Rollstuhls.
